\chapter{Архитектура системы}
\label{ch:architecture}

\section{Общий контекст}
\index{архитектура}

ЦР координирует задания и телеметрию; Регулятор выполняет сертификацию поставки ПО АБУ; взаимодействие компонентов в прототипе осуществляется по \textbf{REST API}. На рисунке~\ref{fig:context} показан контекст развёртывания (PNG генерируется из PlantUML в репозитории).

\begin{figure}[htbp]
  \centering
  \diagram{../diagrams/png/context.png}{0.85\textwidth}
  \caption{Контекст системы: ЦР, АБУ, Регулятор}
  \label{fig:context}
\end{figure}

\section{Цифровой рудник}
\index{ЦР}

Код: \texttt{external\_systems/digital\_mine/}.

\begin{itemize}
  \item регистрация установок (\texttt{rig\_id}, URL API АБУ, опционально \texttt{certificate\_id});
  \item политика \texttt{CR\_CERT\_POLICY}: \texttt{strict} (миссии только при действительном сертификате) или \texttt{permissive};
  \item проксирование создания миссий к API АБУ.
\end{itemize}

\section{АБУ (заготовка)}
\index{ДВБ}

Заготовка АБУ в \texttt{src\_starting\_point/} выполняет миссии и реализует упрощённый «псевдо-ИИ» (эвристики вибрации, режима, флага риска). В версии~1 проверки безопасности не изолированы в отдельной недоверенной зоне --- это \textbf{технический долг} для доработки участниками. Внутренняя структура показана на рис.~\ref{fig:abu-internal}.

\begin{figure}[htbp]
  \centering
  \diagram{../diagrams/png/abu_v1_internal.png}{0.88\textwidth}
  \caption{Внутреннее устройство АБУ (учебный прототип, v1)}
  \label{fig:abu-internal}
\end{figure}

\section{Регулятор}
\index{сертификация}

Регулятор (\texttt{external\_systems/regulator/}) принимает заявку на сертификацию, выполняет прогон \texttt{pytest} с покрытием, отдельный прогон тестов безопасности, оценивает стоимость в зависимости от состава SBOM и ДВБ. Идентификатор \texttt{certificate\_id} в прототипе соответствует хэшу архива сертификационного пакета (см. главу~\ref{ch:certification}).

\section{Сценарий миссии и пайплайн сертификации}

Последовательность взаимодействий при миссии приведена на рис.~\ref{fig:sequence}; этапы подготовки пакета и проверок --- на рис.~\ref{fig:cert-pipe}.

\begin{figure}[htbp]
  \centering
  \diagram{../diagrams/png/sequence_mission.png}{0.92\textwidth}
  \caption{Диаграмма последовательности сценария миссии}
  \label{fig:sequence}
\end{figure}

\begin{figure}[htbp]
  \centering
  \diagram{../diagrams/png/certification_pipeline.png}{0.92\textwidth}
  \caption{Пайплайн сертификации}
  \label{fig:cert-pipe}
\end{figure}

\section{REST: примеры эндпоинтов}

\begin{tabularx}{\textwidth}{@{}l X@{}}
\toprule
\textbf{Система} & \textbf{Примеры путей} \\
\midrule
АБУ & \texttt{/api/v1/health}, \texttt{/api/v1/missions}, \texttt{/api/v1/status} \\
ЦР & \texttt{/api/v1/rigs}, \texttt{/api/v1/missions}, \texttt{/api/v1/security/context} \\
Регулятор & \texttt{/api/v1/certification/requests}, \texttt{/api/v1/certificates/\{id\}}, \texttt{/api/v1/certificates/\{id\}/sga} \\
\bottomrule
\end{tabularx}

\section{Монитор безопасности и учебный пример «Светофор»}
\label{sec:monitor-traffic}
\index{монитор безопасности}
\index{политики безопасности}

\subsection{Шаблон «Монитор (безопасности)».} 

В конкурсном решении предполагается выделение компонента \textbf{монитора безопасности} (в коде --- ориентир \texttt{security\_monitor}) и набора \textbf{политик безопасности} (\texttt{policies}), согласованных с критериями C19--C20 (глава~\ref{ch:evaluation}). 

Монитор активен на \textbf{границах доменов безопасности} АБУ: взаимодействие между недоверенной зоной и зоной, критичной для ЦПБ, не проходит «в~обход», а через проверки, опирающиеся на политики.

\textbf{Изоляция доменов.} Участники должны явно разделить функциональность по зонам доверия: что относится к ДВБ, что может оставаться вне неё, и где проходят контролируемые интерфейсы. Без изоляции компрометация вспомогательного кода распространяется на критичные команды и телеметрию.

\textbf{Контроль по политикам.} Монитор применяет \textbf{политики безопасности} к запросам и ответам между доменами (разрешить, запретить, ограничить параметры, залогировать). Имена и семантика политик должны быть согласованы с моделью угроз и с ЦПБ в SGA.

\textbf{Архитектура и политики.} Архитектура АБУ должна быть согласована не только с политикой информационной безопасности (на уровне системы и эксплуатации), но и с \textbf{теми политиками, которые реально использует монитор}: если объекты политик в коде не покрывают заявленные ЦПБ, формальное наличие монитора не гарантирует нужного поведения при атаках и сбоях.

\subsection{Учебный пример «Светофор».} 

В состав конкурсного задания \textbf{в явном виде} входит ориентир по проектированию изоляции и обмена по правилам --- учебный Jupyter-ноутбук «Светофор», который уже включён в репозиторий: \path{references/traffic_light_demo/cyberimmunity-traffic-lights-example.ipynb}. Краткое руководство по его использованию: \path{docs/traffic_light_demo.md}. Пример иллюстрирует принципы разделения доменов и посредничества монитора; применяется косвенно при оценке C18--C22.

\section{TARA в конкурсном задании}
\label{sec:tara-contest}
\index{TARA}

Методика \textbf{TARA} (анализ угроз и оценка рисков по ISO/SAE~21434) отвечает на вопрос «чему доверять и от чего защищаться».

В рамках конкурсного задания TARA применяется для:
\begin{itemize}
  \item оценки \textbf{требуемых уровней доверия} к компонентам и сценариям эксплуатации;
  \item обоснования \textbf{включения модулей в доверенную вычислительную базу} (ДВБ) и, наоборот, вынесения функций в недоверенную зону при наличии монитора;
  \item трассировки от выявленных угроз к архитектурным мерам (в т.\,ч. к границам доменов и к политикам).
\end{itemize}

Упрощённый разбор угроз и пути атаки (в т.\,ч. для цепочки numpy $\rightarrow$ критичные команды) представлен на рисунке~\ref{fig:tara-numpy} (исходная диаграмма в репозитории: \texttt{docs/diagrams/tara\_attack\_numpy.puml}; растр --- \path{docs/diagrams/png/tara_attack_numpy.png}). Дополнительный разбор --- в \texttt{docs/tara\_abu.md}. Связь с оценкой по критериям --- в главе~\ref{ch:evaluation}.

\begin{figure}[htbp]
  \centering
  \diagram{../diagrams/png/tara_attack_numpy.png}{0.92\textwidth}
  \caption{TARA: путь атаки по цепочке \texttt{numpy} к критичным функциям АБУ}
  \label{fig:tara-numpy}
\end{figure}

\textbf{Вывод.} Если не предпринять дополнительных \textbf{архитектурных и программных} мер (изоляция доменов, вынос тяжёлых вычислений из доверенной зоны, монитор и политики на границах), из данного анализа пути атаки следует необходимость учёта \textbf{numpy} в составе доверенной базы --- фактически \textbf{включения в ДВБ} с типичными для прототипа последствиями: множитель стоимости сертификации для тяжёлых зависимостей в ДВБ и рост \texttt{estimated\_cost}, а через него --- годовой оценки сопровождения на стороне ЦР (см. раздел~\ref{sec:dvb-cost} и формулы~\eqref{eq:cost-heavy}--\eqref{eq:support-annual}).
